\section{What Normalisation Makes of Your Document}
\label{sec:ch17-what-normalization-makes}
\index{router!normalisation pipeline|(}%

The router does not plan the document you sent. It plans what is left of it
after four passes, and since the plan cache is keyed on a hash of that residue,
what those passes throw away is exactly what your clients can vary for free.
Reading them in order is the only way I found to predict which of two documents
would share a plan, and I got two of my own predictions wrong before I did.

\index{router!the four steps, and where validation sits}%
The four calls are in the prehandler, in this order:

\begin{minted}[fontsize=\footnotesize,breaklines]{text}
NormalizeOperation
ValidateOperation
NormalizeVariables
RemapVariables
\end{minted}

Validation sits in the middle of normalisation rather than before it or after
it, which surprised me and has a decent reason: the first pass is what reduces a
multi-operation document to the one operation being run, and validating before
that would mean validating operations nobody asked for. It also means the third
and fourth passes run on a document already known to be valid, so they can
rewrite it without checking anything.

\subsection*{The first pass, and what it is not}

\index{router!static normaliser options}%
The static normaliser is built with four options:

\begin{minted}[fontsize=\footnotesize,breaklines]{go}
func buildNormalizationOptions(enableDefer bool, validateInlineArguments config.ValidateInlineArguments) []astnormalization.Option {
	opts := []astnormalization.Option{
		astnormalization.WithRemoveNotMatchingOperationDefinitions(),
		astnormalization.WithInlineFragmentSpreads(),
		astnormalization.WithRemoveFragmentDefinitions(),
		astnormalization.WithRemoveUnusedVariables(),
	}
\end{minted}

Read that list and then read what is missing from it, because the missing one is
the interesting one. There is no \code{WithExtractVariables()}.
Chapter~\ref{ch:enter-the-router} measured that the planner lifts literal
arguments out into variables, and I assumed on first reading that this pass did
it. It does not. That happens later, and the distinction matters because the two passes
have different jobs and only one of them runs on every request shape.

What the first pass does, in the order the stages are assembled: drop every
operation definition but the one being run; statically resolve a
\code{@skip} or \code{@include} whose argument is a literal boolean, deleting
the selection outright; inline every named-fragment spread, in place, at the
spread's position; remove self-aliasing, which means only \code{foo: foo} and
never a genuinely different alias; splice inline fragments into their parents;
merge sibling inline fragments; drop the fragment definitions that are now
unused; and deduplicate identical fields, keeping the position of the first
occurrence.

\index{router!nothing sorts a selection set}%
Nothing in that list sorts anything. I went looking for a pass that puts a
selection set into a canonical order and there is none, in any of the stages,
which is a fact section~\ref{sec:ch17-what-the-plan-cache-keys-on} cashes in.

\subsection*{Where a literal becomes a variable}

\index{router!variable extraction}%
The lifting happens in the third call, through a different normaliser built for
the purpose. For every argument whose value is a literal, the value is
serialised to JSON, written into the document's variables under a generated
name, and the argument rewritten to point at that name. The names come from
here:

\begin{minted}[fontsize=\footnotesize,breaklines]{go}
const (
	alphabet = `abcdefghijklmnopqrstuvwxyz`
)

func (d *Document) GenerateUnusedVariableDefinitionName(operationDefinition int) []byte {
	var i, k int64

	for i = 1; i < math.MaxInt64; i++ {
		out := make([]byte, i)
		for j := range alphabet {
			for k = 0; k < i; k++ {
				out[k] = alphabet[j]
			}
			_, exists := d.VariableDefinitionByNameAndOperation(operationDefinition, out)
			if !exists {
				return out
			}
		}
	}

	return nil
}
\end{minted}

So \code{a}, then \code{b}, and after \code{z} it goes to \code{aa} and
\code{bb} rather than \code{ab}. The fourth pass renames your own declared
variables by the same scheme, positionally, by first use. Whatever you called
them is gone.

\index{router!the value travels beside the document}%
The half that decides everything downstream is where the literal went. It is not
in the document any more; it is in a JSON side channel that travels alongside.
The document says \code{first: \$a} whether you wrote \code{3} or \code{7}, and
the difference between those two requests lives somewhere the next step never
looks. Figure~\ref{fig:ch17-normalization-to-key} is the whole path with that
branch on it.

\begin{figure}[htbp]
  \centering
  % Four steps the router runs over a client document before it prints and
% hashes the result to a uint64 plan cache key. Referenced from chapter 17.
% Bare tikzpicture; the figure environment, caption and label live at the
% call site.
%
% The point of the picture is what survives to the hash and what does not.
% ValidateOperation sits between the two normalisations and changes nothing
% in the document, which is worth showing because it is surprising.
% NormalizeVariables is where a literal argument value is lifted out into a
% JSON side channel that never reaches the hash, so two requests differing
% only in that value land on the same cache key. The operation name is
% blanked only in the bytes fed to the hash, not in the document itself.
\begin{tikzpicture}[
    font=\small,
    xform/.style={draw, rounded corners=2pt, minimum width=22mm,
                  minimum height=9mm, align=center, font=\scriptsize},
    gate/.style={draw, dashed, rounded corners=2pt, minimum width=22mm,
                 minimum height=9mm, align=center, font=\scriptsize},
    final/.style={draw, thick, rounded corners=2pt, minimum width=26mm,
                  minimum height=9mm, align=center, font=\scriptsize,
                  fill=black!8},
    side/.style={draw, rounded corners=2pt, minimum width=26mm,
                 minimum height=9mm, align=center, font=\scriptsize,
                 fill=black!15},
    flow/.style={-{Stealth[length=2mm]}},
    branch/.style={-{Stealth[length=2mm]}, dashed},
    note/.style={font=\scriptsize, gray, align=center, inner sep=1pt},
    hdr/.style={font=\scriptsize\itshape, align=center}
  ]

  % -- the four steps, in source order -----------------------------------
  \node[hdr] at (4.05,2.5) {four steps over the client document, in source order};

  \node[xform] (n1) at (0.0,1.5) {\code{Normalize}\\\code{Operation}};
  \node[gate]  (n2) at (2.7,1.5) {\code{Validate}\\\code{Operation}};
  \node[xform] (n3) at (5.4,1.5) {\code{Normalize}\\\code{Variables}};
  \node[xform] (n4) at (8.1,1.5) {\code{Remap}\\\code{Variables}};

  \draw[flow] (n1) -- (n2);
  \draw[flow] (n2) -- (n3);
  \draw[flow] (n3) -- (n4);

  \node[note] at (2.7,0.7) {checks the document, changes nothing;\\
    sits between the two normalisations};

  % -- print, then hash -----------------------------------------------------
  \node[xform] (p) at (8.1,-1.6)  {print,\\name blanked};
  \node[final] (h) at (11.0,-1.6) {hash to \code{uint64}\\the plan cache key};

  \draw[flow] (n4) -- (p);
  \draw[flow] (p) -- (h);

  \node[note, anchor=north] at (9.55,-3.1)
    {name blanked only in the hash input;\\
     the document itself still names the operation};

  % -- the literal argument's side channel -----------------------------------
  \node[side] (j) at (5.4,-1.6) {\code{3}\\JSON side channel};

  \draw[branch] (n3) -- (j);

  \node[note, anchor=north] at (5.4,-2.2)
    {\code{first: 3} becomes \code{first: \$a};\\
     same key for \code{first: 3} and \code{first: 7}};

\end{tikzpicture}

  \caption{The four passes, and the one thing that leaves the document on the
  way through. A literal argument is lifted into a generated variable at the
  third step and its value continues beside the document rather than inside it,
  so the hash at the end never sees it. The operation name is removed for the
  hash alone and put straight back, which is why the plan document the router
  will show you still carries a name the cache key ignored.}
  \label{fig:ch17-normalization-to-key}
\end{figure}

\subsection*{The name comes off, and goes back on}

\index{router!the operation name is blanked for the hash}%
The last step before the hash is the one with the comment that explains it:

\begin{minted}[fontsize=\footnotesize,breaklines]{go}
	// Print the operation without the operation name to get the pure normalized form
	// Afterward we can calculate the operation ID that is used as a stable identifier for analytics

	o.kit.normalizedOperation.Reset()
	// store the original name of the operation
	nameRef := o.kit.doc.OperationDefinitions[o.operationDefinitionRef].Name

	staticNameRef := o.kit.doc.Input.AppendInputBytes([]byte(""))
	o.kit.doc.OperationDefinitions[o.operationDefinitionRef].Name = staticNameRef

	err = o.kit.printer.Print(o.kit.doc, o.kit.normalizedOperation)
\end{minted}

The name is saved, replaced with an empty string, the document printed, and that
print is hashed. Then the name goes back and the document is printed a second
time, and \emph{that} second print is the one other things in the router use. So
the operation name is present in the normalised representation and absent from
the plan cache key, and those two facts are going to collide in the next
section in a way I find genuinely awkward.

\medskip
That is what the router makes of a document. What it does with the resulting
number is the question four chapters have been carrying.

\index{router!normalisation pipeline|)}%
